\section{Conclusiones y comentarios}
Se desarrolló el diseño de las fundaciones del edificio, resolviendo los 21 muros mediante zapatas corridas y las 9 columnas mediante zapatas aisladas, a partir de las solicitaciones estáticas y sísmicas del modelo estructural. La totalidad de los 30 elementos verifica las condiciones de presión de contacto, corte, punzonamiento y estabilidad al deslizamiento, conforme a ACI 318-19 y a las solicitaciones sísmicas obtenidas según NCh433 y DS 60.

En cuanto a la presión de contacto, esta resulta la verificación que gobierna el dimensionamiento de la mayoría de las fundaciones. Los cinco muros perimetrales del sector, P1, PN, P13, Plb y P12, se resuelven mediante zapata excéntrica, dado que su base no puede desarrollarse hacia el exterior de la línea de cierro del edificio. En P1, PN, P13 y Plb la excentricidad geométrica resultante reduce la fracción de base comprimida a un rango de 67\% a 68\% y eleva la presión de contacto al orden de 33 a 35~tonf/m$^2$, aún con un factor de seguridad cercano a 2,0 respecto del admisible estático de 70~tonf/m$^2$. P12, en cambio, pese a su misma condición de borde, presenta una presión considerablemente menor, de 11,14~tonf/m$^2$, con la base en compresión completa, dada su menor solicitación. En los muros interiores más solicitados —PE, PJ, PI y PG— la presión sísmica es determinante y el diseño queda ajustado al límite admisible de 93~tonf/m$^2$, con factores de seguridad del orden de 1,00, lo que refleja un dimensionamiento optimizado. En todos los elementos con zapata centrada la zona comprimida supera el 80\% de la base exigido como criterio; el único caso con compresión parcial es el muro PH, que bajo la combinación sísmica más desfavorable alcanza la mayor excentricidad de carga del modelo ($e = 1,10$~m) y conserva aun así un 91,6\% de su base comprimida.

Respecto a la altura de fundación, se adoptó un valor uniforme $H = 1,20$~m para la totalidad de las zapatas, incluidas las de zapata excéntrica. Esta decisión se sustenta en la uniformidad constructiva —evitando sellos de fundación a cotas dispares en las zonas donde las zapatas quedan próximas o en contacto— y en la condición de zapata rígida. La altura útil resultante, $d = 1,15$~m, hace que las verificaciones de corte y punzonamiento queden satisfechas en todos los elementos, con razones de solicitación a resistencia inferiores a 0,23 incluso en los muros perimetrales, por lo que dichas condiciones no controlan el valor de $H$.

En lo que respecta a la estabilidad al deslizamiento, la totalidad de las fundaciones supera el factor de seguridad mínimo de 1,5. Los casos más ajustados corresponden a los cinco muros perimetrales con zapata excéntrica: P1 ($FS_d = 1,50$), P13 y Plb ($FS_d = 1,51$), P12 ($FS_d = 1,51$) y PN ($FS_d = 1,52$), todos prácticamente en el límite. En estos elementos es precisamente la estabilidad al deslizamiento, y no la presión admisible, la verificación que fija el ancho de la base, dado que su holgura en presión resulta menor que la de los muros interiores. El resto de los muros interiores y la totalidad de las columnas presentan factores de seguridad sustancialmente mayores, dada la mayor proporción entre la carga axial estabilizante y el corte solicitante.

En relación a la armadura, en la mayoría de las fundaciones gobierna la armadura mínima por retracción y temperatura, resultando una disposición de $\phi 18 @ 100$~mm. Requieren además una parrilla de refuerzo diseñada por flexión los muros Plb, P3-L, P3-L-J, PL-6, PL-J-11, PI, P5, PJ, PG y PE, correspondientes a los elementos de mayor solicitación de flexión sobre el suelo, para los cuales se adopta una parrilla de $\phi 10 @ 200$~mm.

Finalmente, en relación a las vigas y cadenas de fundación, se observa que varias zapatas quedan próximas o en contacto entre sí, particularmente en el conjunto perimetral P1-PN-P13-Plb-P12 y en el grupo interior conformado por PL-6, PJ2, P5, PJ y PI. En estas zonas se dispone una viga de fundación que vincula los elementos contiguos, controlando los asentamientos diferenciales y aportando arriostramiento ante la acción sísmica. Las zapatas aisladas de columna se conectan mediante cadenas de amarre, que restringen su desplazamiento relativo y absorben las excentricidades transmitidas a la fundación. Su disposición se indica en el plano de fundaciones, sin que su diseño forme parte del alcance de esta entrega.